%
% 6.006 problem set 2 solutions
%
\documentclass[12pt,twoside]{article}
\input{macros-sp20}
\newcommand{\theproblemsetnum}{2}

\usepackage{enumitem}

\title{6.006 Problem Set 2}

\begin{document}

\handout{Problem Set \theproblemsetnum}{}
\setlength{\parindent}{0pt}
\medskip\hrulefill\medskip

{\bf Name:} Nicolas Iskandar

\medskip\hrulefill


\begin{problems}

%%%%%%%%%%%%%%%%%%%%%%%%%%%%%%%%%%%%%%%%%%%%%
\problem  % Problem 1: Solving Recurrences
%%%%%%%%%%%%%%%%%%%%%%%%%%%%%%%%%%%%%%%%%%%%%

\begin{problemparts}

%--- 1(a) T(n) = 4T(n/2) + O(n) ---
\problempart

\textbf{Recursion Tree.}
Build a tree where the root does $O(n)$ work and branches into $4$ children, each on a subproblem
of size $n/2$. At depth $i$ there are $4^i$ nodes, each of size $n/2^i$ and doing $O(n/2^i)$ work,
so the total work at depth $i$ is
\[
  4^i \cdot O\!\left(\frac{n}{2^i}\right) = O(2^i n).
\]
The recursion bottoms out at depth $\log_2 n$, so summing a geometric series that is dominated by
its last term gives
\[
  \sum_{i=0}^{\log_2 n} O(2^i n) = O(n \cdot 2^{\log_2 n}) = O(n^2).
\]
The number of leaves is $4^{\log_2 n} = n^2$, each contributing $\Theta(1)$ work, so the leaves
alone already force $\Omega(n^2)$ work. Combining the two bounds, $T(n) = \Theta(n^2)$.

\textbf{Master Theorem.}
With $a=4$, $b=2$, $f(n) = O(n)$, we compare $f(n)$ against $n^{\log_b a} = n^2$.
Since $n = O(n^{2-\varepsilon})$ for $\varepsilon = 1$, we are in Case~1, giving
\[
  \boxed{T(n) = \Theta(n^2).}
\]

%--- 1(b) T(n) = 3T(n/sqrt(2)) + O(n^4) ---
\problempart

\textbf{Recursion Tree.}
With $a = 3$, $b = \sqrt2$, $f(n) = O(n^4)$: at depth $i$ there are $3^i$ nodes, each of size
$n/(\sqrt2)^i$ and costing $O\!\left((n/(\sqrt2)^i)^4\right) = O(n^4/4^i)$. The total work at
depth $i$ is
\[
  3^i \cdot O\!\left(\frac{n^4}{4^i}\right) = O\!\left(n^4\Big(\frac{3}{4}\Big)^{i}\right).
\]
Because $3/4 < 1$, this is a decreasing geometric series in $i$, so its sum over all
$\log_{\sqrt2} n = 2\log_2 n$ levels is dominated by the root term:
\[
  \sum_{i=0}^{2\log_2 n} O\!\left(n^4\Big(\tfrac{3}{4}\Big)^i\right) = O(n^4).
\]
The root alone costs $\Theta(n^4)$, so this matches the upper bound and $T(n) = O(n^4)$.
For the lower bound, simply note the root by itself costs $\Omega(n^4)$ (dropping all recursive
calls only decreases the total), so $T(n) = \Omega(n^4)$ as well.
Hence $T(n) = \Theta(n^4)$.

\textbf{Master Theorem.}
With $a = 3$, $b = \sqrt2$, we get $n^{\log_b a} = n^{\log_{\sqrt2} 3} = n^{2\log_2 3} \approx n^{3.17}$.
Since $f(n) = \Theta(n^4)$ and $4 > 2\log_2 3$, we check the regularity condition for Case~3:
\[
  a f(n/b) = 3\left(\frac{n}{\sqrt2}\right)^4 = 3\cdot\frac{n^4}{4} = \tfrac{3}{4}n^4 \le c\, n^4
  \quad \text{for } c = \tfrac{3}{4} < 1.
\]
The condition holds, so Case~3 applies and
\[
  \boxed{T(n) = \Theta(n^4).}
\]

%--- 1(c) T(n) = 2T(n/2) + 5n log n ---
\problempart

\textbf{Recursion Tree.}
At depth $i$ there are $2^i$ nodes of size $n/2^i$, each costing $5 \cdot \frac{n}{2^i}\log\frac{n}{2^i}$.
Writing $\log(n/2^i) = \log n - i$, the total work at depth $i$ is
\[
  2^i \cdot 5\frac{n}{2^i}(\log n - i) = 5n(\log n - i).
\]
Reindexing with $j = \log n - i$ (so $j$ runs from $\log n$ down to $0$ as $i$ runs from $0$ to
$\log n$), the total cost over the whole tree is
\[
  \sum_{j=0}^{\log n} 5nj = 5n \cdot \frac{(\log n)(\log n+1)}{2} = \Theta(n\log^2 n).
\]

\textbf{Master Theorem (Case 2, extended form).}
Here $a=b=2$ so $n^{\log_b a} = n$, and $f(n) = 5n\log n = \Theta(n^{\log_b a}\log^k n)$ with
$k=1$. The extended Case 2 gives $T(n) = \Theta(n^{\log_b a}\log^{k+1} n)$, so
\[
  \boxed{T(n) = \Theta(n\log^2 n).}
\]

%--- 1(d) T(n) = T(n-2) + Theta(n) by substitution ---
\problempart

\textbf{Claim:} $T(n) = \Theta(n^2)$.

\textit{Intuition.} Unrolling once gives $T(n) = \Theta(n) + T(n-2)$, and repeating this
$\Theta(n)$ times peels off terms $\Theta(n), \Theta(n-2), \Theta(n-4), \dots, \Theta(1)$, an
arithmetic-like sum of $\Theta(n)$ terms averaging $\Theta(n)$ each, so we expect $\Theta(n^2)$
total. Substitution turns this guess into a proof by induction: we assume the guessed bound holds
for the smaller subproblem $T(n-2)$, substitute it into the recurrence, and check that the
algebra closes for some valid choice of constant.

Let $d>0$ be the constant hidden inside $\Theta(n)$, so concretely
\[
  T(n) = T(n-2) + dn.
\]

\textit{Upper bound: $T(n) = O(n^2)$.}
\textbf{Guess:} $T(n) \le cn^2$ for some constant $c>0$ and all $n \ge n_0$.

\textbf{Inductive hypothesis:} assume the guess already holds for the smaller input $n-2$, i.e.\
$T(n-2) \le c(n-2)^2$.

\textbf{Substitute} this into the recurrence:
\[
  T(n) = T(n-2) + dn \;\le\; c(n-2)^2 + dn.
\]

\textbf{Expand} $c(n-2)^2 = cn^2 - 4cn + 4c$, so
\[
  T(n) \;\le\; cn^2 - 4cn + 4c + dn \;=\; cn^2 + (d-4c)n + 4c.
\]

We want to conclude $T(n) \le cn^2$, which holds exactly when the leftover terms satisfy
\[
  (d-4c)n + 4c \;\le\; 0.
\]

\textbf{Choose $c$:} pick $c = d$, so $d - 4c = -3d < 0$. The leftover expression becomes
$-3dn + 4d$, a downward-sloping line in $n$ that is $\le 0$ once
\[
  n \;\ge\; \tfrac{4}{3}.
\]
So for all $n \ge n_0 = 2$ (checking the base case $T(1)=\Theta(1)$ separately, absorbed by
choosing $c$ large enough), the leftover term is non-positive, giving $T(n) \le cn^2$ as desired.
By induction this holds for all $n \ge n_0$, so $T(n) = O(n^2)$.

\textit{Lower bound: $T(n) = \Omega(n^2)$.}
\textbf{Guess:} $T(n) \ge cn^2$ for some (possibly different, smaller) constant $c>0$.

\textbf{Inductive hypothesis:} assume $T(n-2) \ge c(n-2)^2$.

\textbf{Substitute and expand} exactly as before:
\[
  T(n) = T(n-2) + dn \;\ge\; c(n-2)^2 + dn \;=\; cn^2 + (d-4c)n + 4c.
\]

Now we want the leftover terms $(d-4c)n + 4c$ to be $\ge 0$, so that $T(n) \ge cn^2$ follows
directly.

\textbf{Choose $c$:} pick $c = d/8$, so $d - 4c = d - d/2 = d/2 > 0$. Then the leftover term
$\tfrac{d}{2}n + 4c$ is positive for every $n \ge 0$ automatically, no need to wait for $n$ to
be large. Hence $T(n) \ge cn^2$ for all $n \ge 1$, so $T(n) = \Omega(n^2)$.

\textit{Conclusion.} Combining $T(n)=O(n^2)$ and $T(n)=\Omega(n^2)$,
\[
  \boxed{T(n) = \Theta(n^2).}
\]

\end{problemparts}

\newpage
%%%%%%%%%%%%%%%%%%%%%%%%%%%%%%%%%%%%%%%%%%%%%
\problem  % Problem 2: Sorting Sorts
%%%%%%%%%%%%%%%%%%%%%%%%%%%%%%%%%%%%%%%%%%%%%

\begin{problemparts}

%--- 2(a) ---
\problempart

\textbf{Algorithm: Selection sort.}

The cost of each sort is dominated by how it uses \texttt{get\_at} (cheap, $\Theta(1)$) versus
\texttt{set\_at} (expensive, $\Theta(n\log n)$), so the right choice is whichever algorithm
minimizes total \emph{writes}, not comparisons.

\begin{itemize}
  \item \emph{Selection sort} performs $\Theta(n^2)$ reads to find each minimum but only $n$
        writes total (one placement per pass): cost $\Theta(n^2) + \Theta(n)\cdot\Theta(n\log n)
        = \Theta(n^2\log n)$.
  \item \emph{Insertion sort} can perform up to $\Theta(n^2)$ writes in the worst case (shifting
        elements one at a time), giving $\Theta(n^2)\cdot\Theta(n\log n) = \Theta(n^3\log n)$.
  \item \emph{Merge sort} is not in-place (it needs auxiliary storage to merge), so it is
        disqualified by the problem's requirement.
\end{itemize}

Since writes are the bottleneck resource here, \textbf{selection sort} wins with
$\Theta(n^2\log n)$.

%--- 2(b) ---
\problempart

\textbf{Algorithm: Merge sort.}

Here reads/writes are free but comparisons cost $\Theta(\log n)$ each, so the deciding factor is
the number of comparisons each algorithm performs.

\begin{itemize}
  \item Selection sort and insertion sort each make $\Theta(n^2)$ comparisons in the worst case,
        for a total cost of $\Theta(n^2\log n)$.
  \item Merge sort makes only $\Theta(n\log n)$ comparisons (each merge step compares once per
        output element), for a total cost of $\Theta(n\log n)\cdot\Theta(\log n) = \Theta(n\log^2 n)$.
\end{itemize}

Since $n\log^2 n = o(n^2\log n)$, \textbf{merge sort} is asymptotically best, at
$\Theta(n\log^2 n)$.

%--- 2(c) ---
\problempart

\textbf{Algorithm: Insertion sort.}

The key resource to minimize here is total work on a nearly-sorted input, and insertion sort is
the only one of the three whose running time is sensitive to how far the input already is from
sorted.

Each of the $\log\log n$ adjacent swaps can create at most one new inversion, so the array has at
most $\log \log n$ inversions total after the disruption (it started with zero, being sorted).
Insertion sort's running time is $\Theta(n + I)$ where $I$ is the number of inversions in the
input, since the inner \texttt{while} loop only performs work proportional to the number of
inversions it resolves, plus $\Theta(1)$ per outer iteration. Thus insertion sort runs in
\[
  \Theta(n + \log\log n) = \Theta(n).
\]
Selection sort always costs $\Theta(n^2)$ and merge sort always costs $\Theta(n\log n)$
irrespective of how close to sorted the input already is, so both are worse here.

Therefore, \textbf{insertion sort} is best, at $\Theta(n)$.

\end{problemparts}

\newpage
%%%%%%%%%%%%%%%%%%%%%%%%%%%%%%%%%%%%%%%%%%%%%
\problem  % Problem 3: Friend Finder
%%%%%%%%%%%%%%%%%%%%%%%%%%%%%%%%%%%%%%%%%%%%%

\textbf{Algorithm (bidirectional exponential search).}

Datum sits at kilometer $k$ or kilometer $n-k$ for some unknown $k \ge 0$, and we want a strategy
whose cost scales with $k$ rather than with $n$, since the island's ends are most urgent.

\textbf{Phase 1 (bracketing).}
Probe alternately from the north and south ends at exponentially growing offsets. Concretely, for
$j = 0, 1, 2, \dots$:
\begin{itemize}
  \item Teleport to kilometer $2^j$ (measuring from the north). Query the tracker.
  \item Teleport to kilometer $n - 2^j$ (measuring from the south). Query the tracker.
\end{itemize}
Stop as soon as one of these two probes ``overshoots'' Datum, i.e., the tracker reports Datum is
now on the side of the probe closer to the relevant end. At that point we know Datum lies
between the previous probe on that side (at offset $2^{j-1}$, or $0$ if $j = 0$) and the current
one (at offset $2^j$).

Since Datum is $k$ kilometers from the nearer end, the overshoot happens at
$j^\ast = \lceil \log_2 k \rceil$, and we have made $O(j^\ast) = O(\log k)$ probes (both north
and south probes together, since they interleave one-for-one).

\textbf{Phase 2 (refinement).}
The bracket identified in Phase 1 has width at most $2^{j^\ast} = O(k)$. Run ordinary binary
search inside this bracket, which takes $O(\log(2^{j^\ast})) = O(\log k)$ further probes to pin
down Datum's exact kilometer.

\textbf{Correctness.} Phase 1 always terminates because probes march monotonically toward the
center and Datum is at a fixed finite position, so eventually a probe passes Datum's position.
The resulting bracket is guaranteed (by the tracker's directional guarantee) to contain Datum's
true location, so binary search within it correctly isolates the exact kilometer.

\textbf{Running time.} Phase 1 costs $O(\log k)$ teleports, Phase 2 costs $O(\log k)$ teleports,
for a total of $O(\log k)$, as required. $\blacksquare$

\newpage
%%%%%%%%%%%%%%%%%%%%%%%%%%%%%%%%%%%%%%%%%%%%%
\problem  % Problem 4: MixBookTube Chat
%%%%%%%%%%%%%%%%%%%%%%%%%%%%%%%%%%%%%%%%%%%%%

\textbf{Design goal.}
Every operation is simple except \texttt{ban}: when a viewer is banned, \emph{all} of their
previously sent messages must vanish from the chat, and doing this by scanning the whole message
history would be far too slow. The whole design problem reduces to: \emph{how do we delete one
viewer's scattered messages without searching the entire chat?} The answer is to give each viewer
a private index of pointers directly to their own messages, so a ban only ever touches that
viewer's own messages, never anyone else's.

\textbf{Data structures.}

\begin{enumerate}
  \item A \textbf{lookup structure} $H$ (a balanced BST, e.g.\ an AVL tree, keyed on viewer ID,
        to guarantee worst-case $O(\log n)$ operations) mapping each viewer $v$ to a record
        storing $v$'s banned flag and the head of $v$'s private list $L_v$.
  \item A \textbf{global doubly-linked list} $M$ containing every currently undeleted message, in
        chronological order (newest at one fixed end).
  \item For each viewer $v$, a \textbf{private singly-linked list} $L_v$ that stores, in order,
        \emph{pointers directly into the nodes of $M$} corresponding to $v$'s own messages, not
        copies of the messages, but pointers to their exact location in $M$.
\end{enumerate}

The key idea is layer (3): because $L_v$ points directly at $v$'s nodes inside $M$, deleting all
of $v$'s messages never requires searching $M$, we already know exactly where each one lives.

\textbf{Operations.}

\textbf{build($V$): $O(n\log n)$.}
Insert all $n$ viewer records into $H$ one at a time; each insertion costs $O(\log n)$, for
$O(n\log n)$ total. Initialize $M$ and every $L_v$ as empty in $O(1)$ each.

\textbf{send($v,m$): $O(\log n)$.}
Look up $v$'s record in $H$, $O(\log n)$; if banned, do nothing. Otherwise, create a node for $m$,
splice it onto the newest end of $M$ in $O(1)$, and push a pointer to that same node onto the
front of $L_v$ in $O(1)$. Total: dominated by the lookup, $O(\log n)$.

\textbf{recent($k$): $O(k)$.}
Since $M$ only ever contains currently undeleted messages, in chronological order, simply walk
backward from the newest end of $M$, collecting message contents one at a time ($O(1)$ per step),
stopping after $k$ nodes (or when $M$ is exhausted). Total: $O(k)$.

\textbf{ban($v$): $O(n_v + \log n)$.}
Look up $v$'s record in $H$ ($O(\log n)$) and mark it banned. Walk $L_v$: for each of the $n_v$
pointers stored there, we already know exactly which node of $M$ to remove, so splice it out of
the doubly-linked list $M$ in $O(1)$ using its stored neighbor pointers (no searching needed).
Total: $O(n_v + \log n)$.

\textbf{Correctness.}
\texttt{recent} only ever traverses nodes still linked into $M$, and \texttt{ban} removes exactly
(and immediately) all of a banned viewer's nodes from $M$ the moment they are banned, so a
banned viewer's messages can never appear in a later \texttt{recent} call. All chronological
ordering and per-viewer bookkeeping invariants are preserved by every operation, since \texttt{send}
only ever appends to the newest end of $M$ and the front of $L_v$.

\newpage

%%%%%%%%%%%%%%%%%%%%%%%%%%%%%%%%%%%%%%%%%%%%%
\problem  % Problem 5: Beaver Bookings
%%%%%%%%%%%%%%%%%%%%%%%%%%%%%%%%%%%%%%%%%%%%%
\begin{problemparts}

%--- 5(a) ---
\problempart
\textbf{Algorithm.}
We think of a booking schedule as a piecewise-constant ``staircase'' function
$f(t)$ giving the number of rooms reserved at time $t$; a schedule\\
$(k_1,s_1,t_1),\dots,(k_m,s_m,t_m)$ is just $f$ chopped into its maximal
constant-value pieces (with $f=0$ implicitly outside the listed intervals).
Since $B$ must satisfy $R = R_1 \cup R_2$, and every talk in $R$ came from
exactly one of $R_1,R_2$, the number of overlapping talks at any time $t$ is
additive:
$$f_B(t) = f_{B_1}(t) + f_{B_2}(t).$$
So it suffices to compute this sum, expressed as a booking schedule.

We compute it with a linear two-pointer sweep, analogous to merging two
sorted lists. Maintain pointers $i_1, i_2$ into $B_1, B_2$ (initially $0$)
and a cursor $x$ marking the rightmost time already emitted into the output
(initially $-\infty$, or the minimum possible start time). Repeat until both
pointers reach the end of their lists:
\begin{itemize}[label=\textbullet]
  \item If one list is exhausted, emit the current booking of the other
        list as-is (skipping any gap up to its start time via
        $x \gets \max(x, s)$), and advance that pointer.
  \item Otherwise let $(k_1,s_1,t_1)$ and $(k_2,s_2,t_2)$ be the current
        bookings of $B_1,B_2$. Advance $x$ past any joint gap:
        $x \gets \max(x,\min(s_1,s_2))$. Then:
        \begin{itemize}[label=\textbullet]
          \item If $t_1 \le s_2$ (or symmetrically $t_2\le s_1$), the two
                bookings don't overlap; emit $(k_1, x, t_1)$ and advance
                $i_1$ (resp. emit $(k_2,x,t_2)$, advance $i_2$).
          \item Else if $x < s_2$ (resp.\ $x<s_1$), only $B_1$'s (resp.\
                $B_2$'s) booking is active so far; emit $(k_1,x,s_2)$
                (resp.\ $(k_2,x,s_1)$) \emph{without} advancing a pointer.
          \item Otherwise both bookings are simultaneously active; emit
                $(k_1+k_2,\, x,\, \min(t_1,t_2))$, and advance whichever
                pointer(s) end exactly at $\min(t_1,t_2)$.
        \end{itemize}
\end{itemize}
Finally, do a single linear pass merging any two adjacent emitted pieces
that share the same room count and abut in time (this can occur because a
booking may be emitted across several iterations); this restores the
invariant that consecutive pieces of a valid schedule have distinct room
counts.

\textbf{Correctness.} At every point in the sweep, $x$ marks the boundary
up to which $f_{B_1}+f_{B_2}$ has been fully and correctly emitted as a
sequence of constant pieces: each case above identifies the next time at
which either $f_{B_1}$, $f_{B_2}$, or their sum changes value, and emits the
constant piece up to that point with the correct combined height $k_1$,
$k_2$, or $k_1+k_2$. Since every relevant breakpoint of $B_1$ and $B_2$ is
eventually reached by one of $i_1,i_2$, the full sum $f_{B_1}+f_{B_2}$ is
reconstructed, which is precisely $f_B$ for $B$ satisfying $R_1\cup R_2$.
The cleanup pass only merges pieces that are truly one contiguous
constant-height run split across iterations, so it preserves $f_B$ exactly
while restoring the schedule invariant.

\textbf{Running time.} Each iteration of the sweep does $O(1)$ work and
emits one piece to the output. An iteration advances at least one of
$i_1,i_2$ \emph{except} for the ``only one side active'' case, which does
not advance a pointer but strictly increases $x$ to a fixed boundary
($s_1$ or $s_2$); this case cannot repeat on consecutive iterations without
an intervening pointer advance, so it contributes only $O(1)$ extra
iterations per pointer advance. Hence the sweep runs in
$O(|B_1|+|B_2|) = O(n)$ time, and the final cleanup pass is a single linear
scan, also $O(n)$. Total: $O(n)$.

%--- 5(b) ---
\problempart
\textbf{Algorithm.}
Apply divide and conquer, exactly as in merge sort, using part (a) as the
combine step. Given $R$ with $|R|=n$:
\begin{itemize}
  \item If $n=1$, say $R=\{(s,t)\}$, return the trivial schedule
        $B=((1,s,t))$.
  \item Otherwise split $R$ arbitrarily into two halves $R_1,R_2$ of size
        $\lfloor n/2\rfloor$ and $\lceil n/2\rceil$, recursively compute the
        satisfying schedules $B_1$ for $R_1$ and $B_2$ for $R_2$, and return
        $B = \textsc{Merge-Bookings}(B_1,B_2)$ from part (a).
\end{itemize}

\textbf{Correctness.} By induction on $n$: the base case is immediate. For
the inductive step, $B_1$ and $B_2$ are, by the inductive hypothesis, the
correct satisfying schedules for $R_1$ and $R_2$ respectively. Since
$R = R_1 \cup R_2$ (as the split partitions $R$), part (a) guarantees that
merging $B_1$ and $B_2$ yields the correct satisfying schedule for $R$.

\textbf{Running time.} Note $|B_i| \le |R_i|$ for any request set (a
schedule has at most one piece per request boundary), so the merge at a
node handling $n'$ requests costs $O(n')$ by part (a). This gives the
familiar recurrence
$$T(n) = 2T(n/2) + O(n),$$
which by the Master Theorem (or a recursion-tree argument identical to
merge sort's) solves to
$$T(n) = O(n\log n).$$
\end{problemparts}

\end{problems}

\end{document}